\paragraph{Theory --\label{sec:theory}}
Effective field theories, such as those developed by Fan et al.~\cite{Fan_2010} and Fitzpatrick et al.~\cite{Fitzpatrick:EFT}, provide a more general framework for DM scattering than the usual SI and SD treatment.
We use the non-relativistic effective field theory (NREFT) developed in Ref.~\cite{Fitzpatrick:EFT}, in which the WIMP-nucleon interaction is treated as a four-particle contact interaction
\begin{equation}
    \mathcal{L}_\text{int} = \mathcal{O}\chi\bar{\chi}\bar{N}N,
\end{equation}
where $\chi$ and $N$ are non-relativistic fields for the WIMP and target nucleon, respectively.
Enforcing momentum conservation and Galilean invariance \cite{Fitzpatrick:EFT}  reduces the operators to a basis of four Hermitian quantities
\begin{equation}
i \frac{\vec{q}}{m_N}, \quad \vec{v}^\perp \equiv \vec{v} + \frac{\vec{q}}{2\mu}, \quad \vec{S}_{\chi}, \quad \vec{S}_{N},
\label{eq:basis}
\end{equation}
where $\vec{q}$ is the momentum transferred from the WIMP to the nucleon, $m_N$ is the nucleon mass, $\mu$ is the reduced mass of the WIMP-nucleon system, $\vec{v}$ is the relative velocity between the WIMP and the nucleon, $\vec{v}^\perp$ is the component of the relative velocity  perpendicular to the momentum transfer, $\vec{S}_{\chi}$ is the spin of the WIMP, and $\vec{S}_{N}$ is the spin of the relevant nucleon. 
Linear combinations of these Hermitian quantities up to 2$^{\rm nd}$-order in $\vec{q}$ give a total of fifteen independent and dimensionless NREFT operators, denoted as $\mathcal{O}_i$. The operators can be mapped onto fully relativistic covariant Lagrangians, the nonrelativistic reductions of which we denote $\mathcal{L}_i$, allowing for more complex interactions to be studied~\cite{PandaX:2023_dm_luminescence,LZ:SR1_NREFT_Lagrangian_2024,Anand:MathematicaEFT}.
Equation~\ref{eq:basis} can also be modified to describe inelastic scatters where the WIMP transitions into a more massive state during the interaction, another mechanism that can lead to a suppression of the low energy component of the recoil spectrum~\cite{Barello_2014}. 
Inelastic models arise outside of NREFT, and they can generate interaction rates with significant annual modulation peaking around June 2~\cite{smith_2021:inelasticDM,Tucker-Smith:inelasticDM,Bramante2016:inelasticDM,Graham:Higgsino_2025,Han_1997}.



Here, we perform a search for high-energy nuclear recoils arising from $\mathcal{L}_{1}-\mathcal{L}_{20}$~\cite{Anand:MathematicaEFT}, and for inelastic interactions from $\mathcal{O}_1$ and $\mathcal{O}_4$
%~\cite{Fitzpatrick:EFT} 
with superscripts ($s$) and ($v$) referring to the isoscalar and isovector basis, respectively.
The inelastic operators are selected because they are extensions of SI and SD models.
We use two interactions as illustrative examples, shown in~\autoref{fig:signal_recoil_spectra}. 
The first is the inelastic spin-independent scatter, $\mathcal{O}^{s}_{1}$, which resembles a  Higgsino model~\cite{Graham:Higgsino_2025}.
The second is $\mathcal{L}^{s}_{10}$, which describes a magnetic moment interaction, chosen as it results in a broad and double-peaked spectrum in our region of interest (ROI). 
Predicted xenon nuclear recoil spectra for all dark matter models are calculated with WimPyDD~\cite{wimpydd:Jeong_2022} using one-body nuclear density matrices developed for DMFormfactor-v6~\cite{Anand:MathematicaEFT}, with modifications as described in Ref.~\cite{LZ:SR1_NREFT_2023}, and the Standard Halo Model with parameters recommended in Ref.~\cite{DM_parameters:BAXTER2021_Conventions}. 
Results for all of the probed interactions are in the \textcolor{black}{Data Release}.

\begin{figure}
    \includegraphics[width=0.95\columnwidth]{\plotfolder Fig1_combined_recoils.pdf}
    \caption{Nuclear recoil energy spectra for the two interactions chosen as examples in this Letter.
    \textbf{Top:} Interactions via the inelastic $\mathcal{O}^{s}_{1}$ with splitting between WIMP energy levels ($\delta$) of 0~keV (equivalent to a purely elastic spin-independent interaction), 200~keV, and 300~keV for an impinging WIMP of mass 1000~GeV/c$^2$. 
    \textbf{Bottom:} Elastic interactions of a 50, 200, and 1000~GeV/c$^2$ WIMP via $\mathcal{L}^{s}_{10}$.
    The shaded gray regions, $<$5.4~keV and $>$269.9~keV, indicate where the detection efficiency is below 50\% after all data analysis criteria.
    %The coupling strength of each interaction is assumed to be unity. 
    The coupling strength of each interaction is assumed to be unity 
    %(as in Eq.~21 of Ref.~\cite{Anand:MathematicaEFT}).
    (\emph{i.e.}, $c_i^s = 1/m_\nu^2$ as in Eq. 21 and $d_i^s = 1/m_\nu^2$ as in Eq. 68 of Ref.~\cite{Anand:MathematicaEFT}, where $m_\nu = 246.2$~GeV/c$^2$ is the Higgs vacuum expectation value.). 
    }
    \label{fig:signal_recoil_spectra}
\end{figure}
